\chapter*{Open Logic Project (Açık Mantık Projesi) Hakkında}
\addcontentsline{toc}{chapter}{Open Logic Project (Açık Mantık Projesi) Hakkında}

\textit{Open Logic Text} (Açık Mantık Metni), biçimsel meta-mantık ve biçimsel yöntemleri orta
düzeyden başlayarak ele alan, açık kaynaklı ve ortaklaşa hazırlanan bir ders
kitabıdır. Başka bir deyişle metin, biçimsel mantığa giriş dersini tamamlamış,
matematik ağırlıklı olmayan bir okur kitlesini---özellikle felsefe ve
bilgisayar bilimi öğrencilerini---hedefler; buna karşın matematiksel bakımdan
titizdir.

Metinde yer alan bazı konular henüz bütünüyle işlenmiş olmayabilir; birçok
kesimin de kapsamlı biçimde gözden geçirilmesi gerekmektedir. İleride metni
yeni konuları kapsayacak şekilde genişletmeyi planlıyoruz. Ayrıca sözlük, ileri
okuma listesi, tarihsel notlar, görseller, daha iyi açıklamalar, sonuçların
felsefe, bilgisayar bilimi ve matematik açısından önemini anlatan kesimler ile
daha çok alıştırma ve örnek eklemeyi düşünüyoruz. Bir yanlış bulur ya da bir
öneride bulunmak isterseniz
\href{https://github.com/OpenLogicProject/OpenLogic/wiki/Contributing}{lütfen
proje ekibine bildirin}.

Proje açık kaynak anlayışıyla yürütülmektedir. Metin yalnızca ücretsiz olarak
sunulmakla kalmaz; LaTeX kaynakları da Creative Commons Atıf lisansı altında
yayımlanır. Bu lisans, uygun atıf yapılması koşuluyla çalışmayı indirme,
kullanma, değiştirme, yeniden düzenleme, başka biçimlere dönüştürme ve yeniden
dağıtma hakkını herkese verir. Ek bilgi için Open Logic Project web sitesini
ziyaret edin: \href{http://openlogicproject.org/}{openlogicproject.org}.
